\documentclass[11pt,a4paper]{article}
\usepackage[T1]{fontenc}
\usepackage[utf8]{inputenc}
\usepackage{lmodern}
\usepackage{amsmath,amssymb,amsthm,mathtools}
\usepackage{geometry}
\geometry{margin=2.6cm}
\usepackage{booktabs,array}
\usepackage{enumitem}
\usepackage[dvipsnames]{xcolor}
\usepackage[hidelinks]{hyperref}
\usepackage{xurl}
\hypersetup{pdftitle={No reading avoids the height inequality},pdfauthor={Toshihisa Urahigashi}}
\setlist{itemsep=2pt,topsep=4pt}
\raggedbottom
\emergencystretch=3em

\newcommand{\Z}{\mathbb{Z}}
\newcommand{\Q}{\mathbb{Q}}
\newcommand{\Zp}{\mathbb{Z}_p}
\newcommand{\I}{\mathcal{I}}
\newcommand{\Ocal}{\mathcal{O}}
\newcommand{\Vbad}{\mathbb{V}_{\mathrm{mod}}^{\mathrm{bad}}}
\newcommand{\qq}{\underline{\underline{q}}}
\newcommand{\TT}{\underline{\underline{\Theta}}}
\newcommand{\Cor}{\mathrm{Cor}}
\newcommand{\IVb}{\mathrm{IV}}
\newcommand{\Ht}{\mathrm{H}}
\newcommand{\rV}{\mathrm{V}}
\newcommand{\rL}{\mathrm{L}}
\newcommand{\rLz}{\mathrm{L0}}
\newcommand{\lo}{\mathrm{lo}}
\newcommand{\hi}{\mathrm{hi}}
\newcommand{\leanus}{\textunderscore\allowbreak}
\newcommand{\lean}[1]{{\ttfamily\let\_\leanus #1}}

% labels
\newcommand{\lab}[2]{{\upshape\fcolorbox{#1}{#1!8}{\scriptsize\sffamily\bfseries\textcolor{#1}{#2}}}}
\newcommand{\tagS}{\lab{blue}{S}}
\newcommand{\tagX}{\lab{Purple}{X}}
\newcommand{\tagP}{\lab{ForestGreen}{P}}
\newcommand{\tagR}{\lab{teal}{R}}
\newcommand{\tagN}{\lab{BrickRed}{N}}
\newcommand{\tagH}{\lab{BurntOrange}{H}}
\newcommand{\known}[1]{\par\noindent{\footnotesize\upshape #1}\par\smallskip}
\newcommand{\q}[1]{``#1''}
\newcommand{\srcquote}[2]{\begin{quote}\small #1\par\nopagebreak\hfill{\footnotesize #2}\end{quote}}

\theoremstyle{plain}
\newtheorem{theorem}{Theorem}[section]
\newtheorem{proposition}[theorem]{Proposition}
\newtheorem{corollary}[theorem]{Corollary}
\newtheorem{lemma}[theorem]{Lemma}
\theoremstyle{definition}
\newtheorem{definition}[theorem]{Definition}
\newtheorem{example}[theorem]{Example}
\theoremstyle{remark}
\newtheorem{remark}[theorem]{Remark}

\title{No reading avoids the height inequality:\\ 36 readings of Corollary~3.12\\ and the conservation of the
height load}
\author{Toshihisa Urahigashi\\\small ORCID 0009-0004-0460-6242\\\small\texttt{123026.urahigashi.toshihisa@gmail.com}}
\date{Version 1.1.0, 8 November 2026\\\small\href{https://doi.org/10.5281/zenodo.24039368}{doi:10.5281/zenodo.24039368}}

\begin{document}
\maketitle

\begin{abstract}
Corollary~3.12 of inter-universal Teichm\"uller theory (IUT)~III compares the log-volume $-|\log\qq|$ of the
$q$-pilot object with the log-volume $-|\log\TT|$ of the holomorphic hull of the union of the possible images of the
$\Theta$-pilot object. Theorem~1.10 of IUT~IV bounds $-|\log\TT|$ from above, and the two together give the height
inequality. The value of $-|\log\TT|$ depends on how the statement is read, along five axes: the range of the label
permutations of (Ind1), container or iteration for (Ind3), the group of (Ind2), whether the hull is taken, and the
archimedean container. Combining the axes found so far gives 36 readings that change the value.

We show that in each of the 36 readings, also with a bounded Kummer distortion added, Corollary~3.12 and the upper
bound of IV together give the same height inequality $(k-1)|\log\qq|\le C$. The reason is that $|\log\qq|$, $k$ and
$C$ are computed by one procedure that takes no reading; only the middle term $-|\log\TT|$ depends on the reading.

Moreover, a reading decides only how much of the wall $(k-1)|\log\qq|$ is carried by Corollary~3.12 and how much by
IV. This split is decided by the (Ind1)-axis alone; the other axes change only a bounded quantity that does not involve
the $q$-orders. No reading makes both loads nonpositive. In the family obtained by replacing $q$ with $q^N$, the
statement that holds for every $N$ carries no positive load, and the other carries at least the whole wall.

We prove these statements in Lean~4 for a component model (96 declarations, standard axioms only), and compute the 36
values for seven explicit initial $\Theta$-data. No claim is made about the truth of Corollary~3.12, Theorem~1.10 of
IUT~IV, or the abc conjecture.
\end{abstract}

\tableofcontents

\section{Introduction}

\subsection{Motivation}
The ultimate aim of this work is to understand exactly the algorithm by which IUT computes the height inequality.
Earlier notes determined the value of one step of the computation of IUT~IV that had been left open, the
(Ind2)-term~\cite{Gap}, and put the computation from an initial $\Theta$-datum to the two sides of the height inequality
into one program~\cite{WT}.

That computation contains several choices that the text of the original papers does not fix uniquely; we call them
\emph{readings}. This raises a question: by choosing the readings differently, is there a route that avoids the height
inequality, or one that makes the proof cheaper? The answer of this note is that no reading provides either.

\subsection{The two statements and the readings}\label{ss:two}
Corollary~3.12 compares two log-volumes. One of them is subject to the indeterminacies; the other is not.

\srcquote{\tagS\ \q{Write $-|\log(\TT)|\in\mathbb{R}\bigcup\{+\infty\}$ for the procession-normalized mono-analytic
log-volume [i.e., where the average is taken over $j\in\mathbb{F}_l^{\divideontimes}$ --- cf.\ Remark 3.1.1, (ii),
(iii), (iv); Proposition 3.9, (i), (ii); Theorem 3.11, (i), (a)] of the holomorphic hull [cf.\ Remark 3.9.5, (i)] of
the union of the possible images of a $\Theta$-pilot object [cf.\ Definition 3.8, (i)], relative to the relevant
Kummer isomorphisms [cf.\ Theorem 3.11, (ii)], in the multiradial representation of Theorem 3.11, (i), which we regard
as subject to the indeterminacies (Ind1), (Ind2), (Ind3) described in Theorem 3.11, (i), (ii).}\par
\q{Write $-|\log(\qq)|\in\mathbb{R}$ for the procession-normalized mono-analytic log-volume of the image of a
$q$-pilot object [cf.\ Definition 3.8, (i)], relative to the relevant Kummer isomorphisms [cf.\ Theorem 3.11, (ii)],
in the multiradial representation of Theorem 3.11, (i), which we do not regard as subject to the indeterminacies
(Ind1), (Ind2), (Ind3) described in Theorem 3.11, (i), (ii).}}{\cite{IUT3}, Corollary~3.12}

The proof of Theorem~1.10 of IV bounds $-|\log\TT|$ from above. In doing so it applies Proposition~1.4\,(iii) of IV
with the position of the theta value fixed at the label $j$:

\srcquote{\tagS\ \q{We would like to apply Proposition 1.4, (iii), to the present situation, by taking}\ \dots\
\q{`$i^\dagger$' to be $j\in\mathbb{S}^{\pm}_{j+1}$;}\ \dots\ \q{`$\lambda$' to be `$\mathrm{ord}(-)$' of the
element $\qq^{\,j^2}_{\underline{v}_j}$}\ \dots\ \q{if $\underline{v}_j\in\underline{\mathbb{V}}^{\mathrm{bad}}$}}{\cite{IUT4},
proof of Theorem~1.10, Step~(v)}

This gives the chain
\[-|\log\qq|\ \le\ -|\log\TT|\ \le\ U(D),\]
whose two ends together give the height inequality. How the middle term $-|\log\TT|$ is formed depends on the range of
the possible images, the container, the group, the hull and the archimedean part; such a choice is what we call a
\emph{reading}. For each reading there are two statements, Corollary~3.12 and the upper bound of IV. We show that no
such pair avoids the height inequality, and that a reading only moves the content of the wall between the two
statements.

\subsection{Results}
\begin{itemize}
\item Theorem~\ref{t:cons} (conservation): in every reading, Corollary~3.12 and the upper bound of IV together give the
same height inequality $(k-1)a\le C$, where $a=|\log\qq|$. The sum of the slacks of the two statements equals the slack
of the height inequality.
\item Theorem~\ref{t:split} (the split of the load): with the split point $x$, Corollary~3.12 carries the load
$(x-1)a$ and IV carries the load $(k-x)a$. Their sum is always the wall $(k-1)a$. The two loads are never both
nonpositive, and one of them carries at least half of the wall.
\item Theorem~\ref{t:ind1} (the split point is decided by (Ind1) alone): in the component model the theta part is
$-ka$ for V, increases in the order V, L, L0, and is homogeneous of degree one in the $q$-orders. The theta
part of V equals those of L and L0 exactly when the $q$-orders agree at the places of positive weight above each
prime.
\item Proposition~\ref{p:rest} (the other axes change only the remainder): (Ind2), (Ind3), the hull, the archimedean
part and a bounded Kummer distortion change the value only by a bounded quantity that does not involve the $q$-orders.
\item Theorem~\ref{t:cheap} (the proof cannot be made cheaper): in the family obtained by replacing $q$ with $q^N$, no
reading keeps the two statements together for all large $N$. In a reading in which IV holds for every $N$, the load of
IV is $0$ and Corollary~3.12 carries the whole wall. In a reading in which Corollary~3.12 holds for every $N$, one has
$x\le1$, and IV carries at least the whole wall. The theta part is the slope, in $N$, of the values of the family.
\item Example~\ref{e:two} (two places): above one prime there are a bad place of weight $w\le\frac12$ and a good place
(the situation of Belyi points). In reading L0 one has $x\le1$. In V, IV always holds and Corollary~3.12 fails; in L0
it is the other way round.
\item Section~\ref{s:data} (explicit data): the 36 readings are computed for seven initial $\Theta$-data. All seven
have a prime above which the $q$-orders differ, and $x_\rL$ and $x_\rLz$ are strictly less than $k$. For $D_{211}$,
$D_{109}$ and $D_{15}$ one has $x_\rLz=\frac78<1$.
\end{itemize}
In one sentence: a reading decides only which of the two statements carries the content of the wall; the total does
not depend on the reading.

\subsection{Relation to earlier work}\label{ss:prior}
\begin{itemize}
\item \tagP\ Theorem~14.3\,(ii) of the author's educational note~\cite{II} shows, in a model that treats the (Ind1)-axis
alone, that in every reading Corollary~3.12 and IV together imply the height bound (an implication). There, what (Ind2)
and (Ind3) add to the hull is represented by one number $E$. This note extends that statement to five axes and the
Kummer distortion, turns the implication into an identity (the sum of the loads), and adds the theorems that the split
point is decided by (Ind1) alone, with the condition for equality, and the uniformity theorem for the family $q^N$.
\item \tagX\ Earlier, the public repository Takkun-kohinata/IUT\_LEAN~\cite{TK} (GitHub, 15 July 2026) proved in Lean,
for one reading, that the containment of Step~(v) of IV, with $\lambda=\mathrm{ord}(q^{j^2})$, together with
Corollary~3.12 gives the height inequality $(\ell(\ell+1)/12-1)\,a\le L$, with $a=|\log\qq|$ as here and $L\ge0$ its
ledger of IV, Propositions~1.1--1.4 (\lean{thetaHull\_forces\_height}); that conversely this inequality, together with
the condition $C_\Theta\ge-1$ of IV, Theorem~1.10, gives a model with that containment
(\lean{thetaHullTranscription\_iff}); and that with $L=0$ the containment cannot hold (\lean{no\_lossless\_theta\_hull}),
all in \lean{IutLean/Multiradial/ThetaHullDichotomy.lean}. In the terms of this note this is the case $x=1$, in which
Corollary~3.12 carries no load and IV carries the whole wall.
\item \tagX\ The repository lana-agents/iut~\cite{lana} (``lana'' below), a public Lean formalization, distinguishes
two readings of (Ind1), the literal reading and the reading of IV, and states that the literal reading does not suffice
for the ABC argument. It also proves that taking the hull for each (Ind1)-image separately gives the value of the
reading of IV (\lean{rhsDataOn\_singleton} in \lean{Iut/Concrete/Ind1Invariance.lean}).
\end{itemize}
\srcquote{\tagX\ \q{Under the literal union over the labels, the holomorphic hull at a tuple \lean{c} is governed by
\lean{min\_l ord\_p(q\_\{c(l)\})} rather than by the average of \lean{ord\_p(q\_v)} over the places that IUT IV,
Step (v), obtains via Proposition 1.7. So the literal reading does not suffice for the Theorem 1.10 / ABC argument: the
points produced by the Belyi reduction have good and bad places above the same prime, where
\lean{min\_l ord\_p(q\_\{c(l)\}) = 0}.}}{\cite{lana}, \lean{Iut/Concrete/Literal.lean}, module docstring}
\begin{itemize}
\item \tagX\ In the same repository, \lean{Iut/WallEquiv/} collects formulations of the wall---the word of the interim
report of Project LANA~\cite[\S0.2]{LANA}---and proves equivalences between them in its models, and
\lean{NonDerivable.lean} shows, for one reading, that Step~(xi) does not follow from the formalized content of
Theorem~3.11 alone. lana itself states that these equivalences hold only in the models.
\end{itemize}
\srcquote{\tagX\ \q{\textbf{Only in the models} (identified-copies output law, zero distortion --- e.g.\ unit-induced
Kummer identifications, \lean{Iut/LogLink}): 1~$\Leftrightarrow$~2~$\Leftrightarrow$~3~$\Leftrightarrow$~5~$\Leftrightarrow$~6~$\Leftrightarrow$~10.
A formulation different from (SS) needs an $\eta$-algorithm whose output law is not \lean{LGP} acting on an
(Ind3)-choice, or a transport with nonzero distortion (a Kummer identification that is not unit-induced).}\par
\q{\textbf{The missing input is ABC-strength}: \lean{classicalABC\_of\_forall\_xi} (Step (xi) for all genuine data gives
\lean{ClassicalABC}, via \lean{classicalABC\_of\_thm311}).}}{\cite{lana}, \lean{Iut/WallEquiv/README.md}}
\begin{itemize}
\item \tagX\ lana's file \lean{Iut/Abc/P24Table.lean} (added on 14 October 2026) checks by \lean{decide} a table of the
ingredients that a proof of the Szpiro inequality must use (from a research note \lean{P2P9.md}, \S2.4, which is not
in the repository). A ``reading'' there is a way of deciding the open entries of the table, not a reading of
Corollary~3.12. For proofs through the IUT package (structure C6), the ingredient VA is added to every minimal hitting
set, in the baseline reading of the table.
\end{itemize}
\srcquote{\tagX\ \q{\lean{VA}: the IUT coupling \lean{V $\ge$ (s - 1) a}.}\par
\q{With \lean{C6} (proofs through the IUT package), \lean{VA} is added to each minimal hitting set.}}{\cite{lana},
\lean{Iut/Abc/P24Table.lean}}
\begin{itemize}
\item \tagX\ The shape of $\Ht$ below, $k\,a\le a+C$, is that of the statement in Mochizuki's report on formalization
that the multiradial algorithm leaves the height unaffected, up to a mild error, by multiplication by
$N$~\cite[\S3]{Form}. The Lean port of the toolkit motizukipdf of D.~Richards~\cite{Richards} (October 2026) writes this
as $N\cdot h\le h+C\Rightarrow h\le C/(N-1)$ (\lean{height\_bounded\_of\_dilation}) and proves, from the comparison in
the common container taken as a hypothesis, the bound $|\log\qq|\le H/(N-1)$ with $N=\ell(\ell+1)/12$ and $H$ its
estimate of the hull term (\lean{corThreeTwelve\_of\_thmThreeElevenFive}).
\item \tagH\ Thus the fact that proofs through the IUT package use the wall itself is already in Takkun-kohinata/IUT\_LEAN
and in lana, for one reading (\lean{ThetaHullDichotomy.lean}, \lean{NonDerivable.lean} and VA above). What this note adds
is the form over the whole family of readings. lana does
not state that in every reading the load only moves between the two statements while its total stays the same. For
the literal reading lana states that it does not suffice, but not that the load moves to IV (checked at commit
\texttt{553076a}, the head of the repository on 5 November 2026).
\item \tagX\ The correspondence with the dilemma of Scholze and Stix~\cite{SS} (blurring too much gives an empty
statement, blurring too little a false one) is discussed in Section~\ref{ss:types}.
\item \tagX\ Joshi holds that the Tate parameter changes with (Ind1), (Ind2) and (Ind3)~\cite{Joshi}. This denies the
premise of this note that the left end takes no reading (Section~\ref{ss:change}).
\item \tagP\ The author's note~\cite{W311} uses the same family $q\to q^N$ in its models of the log-volume content of the
results cited in the proof of Corollary~3.12 (its Theorems~14, 15 and~22).
\item \tagN\ What we consider new, to the author's knowledge: the identity of the loads over the 36 readings and the
Kummer distortion; the fact that the split point is decided by (Ind1) alone, the other four axes changing only the
remainder; the uniformity theorem for the family $q^N$ (a statement that holds for every $N$ carries no positive load);
and the computation of the 36 readings on explicit data. The condition for equality is not new in itself: lana's
\lean{Literal.lean} states, from revision \texttt{2b2190d} (11 October 2026), that the regions of its two readings
coincide when the $q$-orders agree at the places above each prime, and that two bad places above $p$ with different
$q$-orders already break this; and \cite[Proposition~13.2]{II} gives the same condition for V and L0 in a model of the
(Ind1)-axis alone. What is added here is the same condition for L, and its validity in all 36 readings. That the wall is
needed is not new (Takkun-kohinata/IUT\_LEAN and lana, above).
\end{itemize}

\subsection{What is not claimed}
\begin{itemize}
\item It is not claimed that the 36 readings are all the readings of the original papers. They are the combinations of
the axes found so far.
\item That each axis changes only a bounded quantity (Proposition~\ref{p:rest}) rests on a paper argument from the
formula for each axis and on checks by the programs. Lean uses it as a hypothesis.
\item The family obtained by replacing $q$ with $q^N$ is a model, not a family of actual curves.
\item We do not decide which reading the original papers intend. Nothing is claimed about the truth of
Corollary~3.12, Theorem~1.10 of IUT~IV, or the abc conjecture.
\end{itemize}

\subsection{Labels}
\tagS\ a statement of the original papers of IUT or of Mochizuki's documents. \tagX\ a statement of others (lana,
Scholze--Stix and others). \tagP\ a statement proved here; \tagP\ alone marks an elementary fact. \tagR\ known
mathematics, applied here in a new way; the known mathematics is named. \tagN\ a result new in this note, to the
author's knowledge. \tagH\ a reading or interpretation, not a proof.

\section{Setting}\label{s:setting}

\subsection{The two statements and the height inequality}
For an initial $\Theta$-datum $D$ (with a prime $\ell\ge5$ and $\ell^*=(\ell-1)/2$) put:
\begin{itemize}
\item $a(D)=|\log\qq|>0$: the negative of the log-volume of the $q$-pilot object;
\item $k=\dfrac1{\ell^*}\displaystyle\sum_{j=1}^{\ell^*}j^2=\dfrac{\ell(\ell+1)}{12}$: the average of $j^2$ over the
labels; $k>1$ for $\ell\ge5$;
\item $U(D)$: the upper bound of $-|\log\TT|$ given by Steps (v)--(vii) of the proof of Theorem~1.10 of IV. As in the
original, the contributions of the primes $2$, $3$, $5$ and $\ell$ are bounded by Step~(ii) and (R4). It is computed by the program
\texttt{iut\_height.py} of this note, which implements Steps (iv)--(viii) as written;
\item $C(D)=U(D)+k\,a(D)$: the upper bound of IV without the term $-k\,a$ of the theta values.
\end{itemize}

\begin{definition}[The two statements and the height inequality]\label{d:two}
Write $M_R(D)$ for the value of $-|\log\TT|$ in a reading $R$, and put
\[\begin{aligned}\Cor_R(D)&:\ -a(D)\le M_R(D),\\ \IVb_R(D)&:\ M_R(D)\le U(D)=-k\,a(D)+C(D),\\
\Ht(D)&:\ (k-1)\,a(D)\le C(D).\end{aligned}\]
$\Ht(D)$ contains no reading.
\end{definition}

\tagP\ The height inequality $\Ht$ is the inequality $H/6\le F_\chi U$ of Proposition~3 of the author's fifth
report~\cite{Fifth}, written in other variables: there $H$ is the degree of the divisor of the $q$-parameters, so that
$a=H/(2\ell)$, and $\chi=12/(\ell(\ell+1))=1/k$, $F_\chi=1/(1-\chi)=k/(k-1)$ and $C=\frac{\ell+1}4U$ (the $U$ of that
report is not the $U(D)$ of this note).

\tagH\ We treat the upper bound $U(D)$ of IV as one number that does not depend on the reading. What the text of IV
computes is this one number, and when the reading is changed, the question is whether that number bounds the hull in
the reading $R$ (the statement $\IVb_R$ of Definition~\ref{d:two}). Recomputing an upper bound for each reading is
treated in Theorem~\ref{t:split}\,(v).

\tagP\ The quantities $a$, $k$, $U$ and $C$ are computed by one procedure that takes no reading ($a$ and the side of
Corollary~3.12 by the program of~\cite{WT}, $U$ by \texttt{iut\_height.py}). That the left end is not subject to the
reading is the wording of Corollary~3.12 itself (Section~\ref{ss:two}), and the procedure for the right end fixes the
position $i^\dagger=j$ in Step~(v). Moreover, we compared the procedure with an elementary computation written from the
statement of Theorem~1.10 alone, without the machinery of IUT (\texttt{elementary\_height.py}, which does not use the
program of~\cite{WT}; both programs are the author's). For all seven data the places counted and the four numbers at
the ends ($\log\mathfrak q$, $\log\mathfrak f$, $\log\mathfrak d$ and the right-hand side of the theorem) agree
(\texttt{compare\_elementary\_iut.py}).

\subsection{Five axes and 36 readings}\label{ss:axes}
\begin{center}\small
\begin{tabular}{@{}>{\raggedright\arraybackslash}p{2.9cm}>{\raggedright\arraybackslash}p{6.1cm}>{\raggedright\arraybackslash}p{5.9cm}@{}}
\toprule
Axis & Choices & Source\\
\midrule
(Ind1): the factors on which the theta value of label $j$ may sit &
V: the factor $j$ only; L: $1,\dots,j$ (permutations fixing $0$); L0: $0,\dots,j$ (all permutations) &
\tagS\ V is Step~(v) of IV. \tagX\ V and L0 are lana's reading of IV and its literal reading. L is (L) of the
catalogue~\cite{Cat}. \tagP\ L0 also moves the label $0$: fixing $\mathcal D_0$ does not fix the label $0$ inside the
larger capsules~\cite[\S6]{Fourth}\\
\addlinespace
(Ind3) &
container: the whole tensor product $\bigotimes\I$ of log-shells as support; iteration: only units and the pure
tensors of the actual images of the same number of iterated log-links &
\tagS\ The container is IV, Proposition~1.4\,(iii) and Step~(v). \tagH\ The iteration is our reading\\
\addlinespace
(Ind2) &
units: $\Zp^\times$ and the field automorphisms; lattice: all $\Zp$-linear automorphisms preserving the log-shell of
each factor &
\tagX\ The lattice group is lana's \lean{tensorAut}. \tagP\ The group of IV, Proposition~1.2 gives the same
value~\cite[\S6.2]{Gap}. \tagP\ At a finite extension of $\Q_p$ the group Ism of IUT~III, Theorem~3.11\,(i) is
$\Zp^\times$ (by \cite[Thm.~F, Thm.~7.4]{MT}), so (Ind2) acts by unit scalars~\cite[\S6, Proposition~9]{Fourth}\\
\addlinespace
Hull &
yes: the union is replaced by its holomorphic hull; no: the union itself &
\tagS\ Corollary~3.12 takes the hull. \tagX\ No hull is lana's \lean{SSIneqNoHull} (``not IUT's container'', in its
docstring)\\
\addlinespace
Archimedean part &
small: $\pi^{j+1}\cdot B_I$; large: the unit ball of the tensor-product metric &
\tagS\ Small is Step~(vii) of IV. \tagX\ Large is what lana keeps, following IUT~III, Proposition~3.2\,(ii)
(\lean{VcI.md}, \S2(e))\\
\bottomrule
\end{tabular}
\end{center}

In the iteration reading without the hull, the possible images form a union of pure tensors, of volume $0$
(log-volume $-\infty$). Leaving out these combinations, there are $3\times(2\times2+2)\times2=36$ readings (Lean:
\lean{card\_admissible}).

\srcquote{\tagS\ \q{Then it follows from Proposition 1.5, (iii), (iv), that $\pi^{j+1}\cdot B_I$ serves as a container
for the `union of possible images of a $\Theta$-pilot object' discussed in Step (iv).}}{\cite{IUT4}, proof of
Theorem~1.10, Step~(vii)}

\tagH\ In addition, the Kummer distortion $\delta$ (the difference of log-volumes carried by a Kummer identification
that is not induced by units; the distortion in lana's \lean{WallOn}) is treated separately, as a continuous
parameter; $\delta$ is added to $M_R$. Readings that do not change the value but decide whether the comparison applies
at all, such as (M$_{\mathrm b}$) and (M$_{\mathrm g}$) of the catalogue~\cite{Cat}, are not counted as axes.

\subsection{The component model}\label{ss:model}
We set up the model in which the Lean proofs are carried out. Take a finite set $V$ of places above one prime $p$,
weights $w_v\ge0$ with $\sum_vw_v=1$, and $q$-orders $Q_v\ge0$ (in units of $\log p$; zero except at bad places). For
$j=1,\dots,\ell^*$, a component for the label $j$ is a tuple $t=(t_0,\dots,t_j)\in V^{j+1}$ of weight
$w(t)=\prod_iw_{t_i}$. Write $s_i(t)=j^2Q_{t_i}$ for the order of the theta value in the factor $i$.

Let the sets of factors allowed by the (Ind1)-axis be $P_\rV(j)=\{j\}$, $P_\rL(j)=\{1,\dots,j\}$ and
$P_\rLz(j)=\{0,\dots,j\}$ (Lean: \lean{pos}). The hull of the images that put the theta value on these factors is
decided by the largest image. The theta part on a component is
\[h_r(t)=\max_{i\in P_r(j)}\bigl(-s_i(t)\bigr)=-\min_{i\in P_r(j)}s_i(t)\qquad(r\in\{\rV,\rL,\rLz\})\]
(\lean{thetaHull}), and its procession-normalized value is
\[\Theta_r=\frac1{\ell^*}\sum_{j=1}^{\ell^*}\sum_{t\in V^{j+1}}w(t)\,h_r(t)\]
(\lean{thetaPart}). The degree of the $q$-pilot is $a=\sum_vw_vQ_v$ (\lean{qdeg}). With several primes, the values for
each prime are multiplied by $\log p$ and added (\lean{thetaSum}, \lean{qdegSum}).

\tagR\ Known mathematics: \cite[Setup~12.1 and Propositions 13.1--13.3]{II}, where the readings are given by sets of
permutations. The hulls of V and L0 are those of $S_{\mathrm b}$ (the permutations fixing $j$) and $S_{\mathrm a}$ (all
permutations).

\section{Main results}\label{s:main}

\begin{definition}[Theta part, remainder, split point, loads]\label{d:split}
Write the value of a reading $R$ as its theta part plus a remainder, $M_R=\Theta_{r(R)}+B_R$, where $r(R)$ is the
(Ind1)-axis of $R$. The theta part is the part proportional to the $q$-orders; the remainder $B_R$ collects the
container, the unit parts, the rounding of the lattice, the defect of the hull, the archimedean part and a bounded
Kummer distortion. With the split point $x_r=-\Theta_r/a$ put
\[\begin{aligned}\text{load of Corollary 3.12}&=(x_r-1)\,a, & \text{load of IV}&=(k-x_r)\,a,\\
\text{room of Corollary 3.12}&=B_R, & \text{room of IV}&=C-B_R.\end{aligned}\]
\end{definition}

\tagH\ For one datum alone this split is only bookkeeping, since every $M_R$ can be written as a theta part plus a
remainder. What decides the split uniquely is the slope in the family obtained by replacing $q$ with $q^N$
(Theorem~\ref{t:cheap}\,(iv)).

\begin{theorem}[Conservation]\label{t:cons}\tagP
\known{Elementary. Lean: \lean{Datum.conservation}, \lean{Datum.slack\_sum}}
In every reading $R$, $\Cor_R$ and $\IVb_R$ together imply $\Ht$. Moreover
\[\bigl(M_R+a\bigr)+\bigl(U-M_R\bigr)=C-(k-1)\,a,\]
so the sum of the slacks of the two statements equals the slack of the height inequality.
\end{theorem}

\begin{theorem}[The split of the load]\label{t:split}\tagP
\known{Elementary. Lean: \lean{Datum.cor\_iff}, \lean{Datum.iv\_iff}, \lean{Datum.load\_sum}, \lean{Datum.x\_le\_k},
\lean{Datum.loadIV\_nonneg}, \lean{Datum.loadIV\_eq\_zero\_iff}, \lean{Datum.half\_load\_le\_max},
\lean{Datum.not\_both\_loads\_nonpos}, \lean{Datum.adapted\_bound}}
Assume that the theta part satisfies $\Theta_r\ge-k\,a$ (in the model this holds by Theorem~\ref{t:ind1}\,(ii)). Then:
\begin{enumerate}[label=(\roman*)]
\item $\Cor_R\iff(x_r-1)a\le B_R$, and $\IVb_R\iff(k-x_r)a\le C-B_R$.
\item The two loads add up to $(k-1)a$, and the two rooms add up to $C$. $\Ht$ is the sum of the two inequalities.
\item $x_r\le k$, and the load of IV is nonnegative. It vanishes only when $\Theta_r=-k\,a$.
\item The larger of the two loads is at least $(k-1)a/2$. In particular the two loads are never both nonpositive.
\item If the upper bound of IV is replaced by the bound $M_R\le\Theta_r+\hi$ that holds automatically in the reading
$R$ (where $B_R\le\hi$), the two statements together give only $(x_r-1)a\le\hi$. Its coefficient $x_r-1$ of $a$ is
smaller than the coefficient $k-1$ of the wall by exactly the load of IV, $k-x_r$.
\end{enumerate}
\end{theorem}

\begin{theorem}[The split point is decided by (Ind1) alone]\label{t:ind1}\tagR
\known{Known mathematics: \cite[Propositions 13.1--13.3]{II} (V and L0); for one label position per capsule, lana's
\lean{rhsDataOn\_singleton}; the condition for equality of lana's two readings is stated in lana's \lean{Literal.lean}
(Section~\ref{ss:prior}). Lean: \lean{thetaPart\_V},
\lean{thetaPart\_mono}, \lean{thetaPart\_ge}, \lean{thetaPart\_nonpos}, \lean{thetaPart\_smul},
\lean{thetaPart\_eq\_V\_of\_const}, \lean{thetaPart\_V\_lt\_L0}, \lean{thetaPart\_V\_lt\_L}, \lean{thetaSum\_V},
\lean{thetaSum\_ge}, \lean{thetaSum\_smul}}
In the model of Section~\ref{ss:model}:
\begin{enumerate}[label=(\roman*)]
\item $\Theta_\rV=-k\,a$. That is, $x_\rV=k$, and in reading V the load of IV is $0$.
\item $\Theta_\rV\le\Theta_\rL\le\Theta_\rLz\le0$. That is, $k=x_\rV\ge x_\rL\ge x_\rLz\ge0$.
\item If the $q$-orders are multiplied by $N\ge0$, then $\Theta_r$ and $a$ are multiplied by $N$, and $x_r$ does not
change.
\item If the $q$-orders are equal at all places of positive weight, then $\Theta_\rL=\Theta_\rLz=\Theta_\rV$.
\item If two places of positive weight have different $q$-orders, then $\Theta_\rV<\Theta_\rLz$. If $\ell^*\ge2$, then
also $\Theta_\rV<\Theta_\rL$ (for $j=1$, L is the same as V).
\end{enumerate}
For sums over several primes (with coefficients $\log p\ge0$), (i)--(iii) hold as well.
\end{theorem}

\begin{proposition}[The other axes change only the remainder]\label{p:rest}\tagP
\known{Elementary, given the formula for each axis; the sources of the formulas are given below. In Lean a
hypothesis (\lean{Family.rest\_lo}, \lean{Family.rest\_hi}).}
In units of $\log p$, let $\rho_u$ be the outer radius of the log-shell, $\rho'_u$ a lower bound for its inner radius,
$\mu_u=1+\delta_p-\rho_u$ with
$\delta_p=[p=2]$, $n=j+1$, and $\kappa(X)=\max\{k:X\subseteq p^k\I\}$. The values on a component are as follows.
\begin{itemize}
\item Container, units, hull: $\sum_u\rho_u-\min_{i\in P}s_i$.
\item Container, lattice, hull: $\sum_u\rho_u-\min_{i\in P}\kappa(\theta_i\I)$, where
$s_i-g\le\kappa(\theta_i\I)\le s_i$ with a local constant $g$ ($g<1$ at small places).
\item Iteration, units: $-\bigl(\min_{i\in P}s_i+\min(0,\sum_u\mu_u)\bigr)$.
\item Iteration, lattice: $\sum_u\rho_u-\min_{i\in P}\min(A_i,B_i)$, where
$A_i=\kappa(\theta_i\Ocal)+\sum_{u\ne i}\kappa(\Ocal_u)$ and $B_i=n(1+\delta_p)+\kappa(\theta_i\I)$.
\item No hull (container): at least the value with the hull minus the shell defect, and at most the value with the
hull; the shell defect is nonnegative and at most $d_I+\sum_{I^*}(3+\log e)/\log p$.
\item Archimedean part: the hull of the large ball exceeds the small container by $\frac j2\log2$ for the label $j$
(by $\frac{\ell^*+1}4\log2$ on average over the labels). Without the hull, the large ball is smaller by a further
$\log(N!)/(2N)$, with $N=2^j$.
\end{itemize}
Each formula is $-\min_{i\in P}s_i$ plus a quantity bounded above and below by local constants that do not involve the
$q$-orders ($\rho$, $\rho'$, the constant $a+b$ of IV, Proposition~1.2, $\mu$, $n(1+\delta_p)$ and $3+\log e$). Hence
the theta part is decided by the (Ind1)-axis alone, the other axes change only the remainder $B_R$, and that change
does not involve the $q$-orders. A bounded Kummer distortion also goes into the remainder. A distortion $d\cdot a$
proportional to $a$ lowers the split point by $d$ (and moves load to IV if $d>0$).
\end{proposition}

Sources of the formulas. \tagX\ The value $\sum\rho_u\log p$ of the hull of the container is lana's
\lean{componentVol\_leastHull\_tensorLogShell\_eq}. \tagR\ The bracketing of $\kappa$ for the lattice group is
Lemma~4.1 and Theorem~3.3 of~\cite{Gap}. \tagH\ The two formulas for the iteration reading are our computation in our
reading. \tagX\ The bound for the shell defect and the values of the archimedean part are lana's (\lean{VcI.lean},
\lean{VcI.md}, \S2(e)). \tagX\ Zeng~\cite{Zeng} gives in closed form the index of the integral structure
$\bigotimes^n\Ocal_K$ in the product of the rings of integers of the components of $\bigotimes^nK$, for a Galois
extension $K/\mathbb{Q}_p$, the index that the hull introduces; it depends on the discriminant of $K/\mathbb{Q}_p$ and
on $n$, not on the $q$-orders. Bounded here means: not involving the $q$-orders; the remainder does depend on $\ell$
and on the local fields.

\begin{theorem}[The proof cannot be made cheaper]\label{t:cheap}\tagN
\known{Lean: \lean{Family.no\_uniform\_reading}, \lean{Family.iv\_uniform}, \lean{Family.cor\_uniform},
\lean{Family.tendsto\_slope}, \lean{modelFamily\_scaled}}
Let $D^N$ be the datum with $q$ replaced by $q^N$. Assume that $a$ and the theta parts are multiplied by $N$, that
$\Theta_r\ge-k\,a$ for every $r$ (in the model this holds by Theorem~\ref{t:ind1}\,(ii)), that $k$ and $C$ do not
change, and that the remainders lie in intervals $[\lo_R,\hi_R]$ that do not depend on $N$ (Theorem~\ref{t:ind1}\,(iii)
and Proposition~\ref{p:rest}). Then in every reading $R$:
\begin{enumerate}[label=(\roman*)]
\item There is $N_0$ such that $\Cor_R(D^N)$ and $\IVb_R(D^N)$ do not hold together for $N\ge N_0$.
\item If $\IVb_R(D^N)$ holds for every $N\ge1$, then $\Theta_r=-k\,a$ (the load of IV is $0$), and
\[\Cor_R(D^N)\iff(k-1)\,Na\le B_R(D^N):\]
Corollary~3.12 carries the whole wall.
\item If $\Cor_R(D^N)$ holds for every $N\ge1$, then $\Theta_r\ge-a$ (that is, $x_r\le1$), and $\IVb_R(D^N)$ implies
$(k-1)\,Na\le C-B_R(D^N)$: IV carries at least the whole wall.
\item $M_R(D^N)/N\to\Theta_r$ as $N\to\infty$. The family decides how the value splits into theta part and remainder.
\end{enumerate}
\end{theorem}

\begin{example}[Two places: the situation of Belyi points]\label{e:two}\tagP
\known{Known mathematics: \cite[Proposition 13.3 and Theorem 14.3\,(iv), (v)]{II}. Lean: \lean{thetaPart\_L0\_two},
\lean{twoPlace\_x\_L0\_le\_one}, \lean{belyi\_cor\_L0}, \lean{belyi\_iv\_L0\_fails}, \lean{belyi\_iv\_V},
\lean{belyi\_cor\_V\_fails}}
Above one prime there are a bad place of weight $w$ and $q$-order $Q_{\mathrm b}>0$ and a good place of weight $1-w$.
Then $x_\rLz=A(w)=\frac1{\ell^*}\sum_{j=1}^{\ell^*}j^2w^j$, and $x_\rLz\le1$ if $w\le\frac12$.

Take the family with $w=\frac12$, $Q_{\mathrm b}=1$ and remainder $0$. In reading V, if $C\ge0$, IV holds for every
$N$ and Corollary~3.12 fails for every $N$. In reading L0, Corollary~3.12 holds for every $N$ and IV fails for large
$N$. The two readings put the load on opposite sides.
\end{example}

\section{Proofs}\label{s:proofs}

\begin{proof}[Proof of Theorem~\ref{t:cons}]
Adding $-a\le M_R$ and $M_R\le-k\,a+C$ gives $-a\le-k\,a+C$, that is, $(k-1)a\le C$. The identity is obtained by
expanding.
\end{proof}

\begin{proof}[Proof of Theorem~\ref{t:split}]
Substituting $\Theta_r=-x_ra$ gives (i). (ii) is the sum of the two inequalities of (i). (iii) follows from
$-k\,a\le\Theta_r=-x_ra$ and $a>0$. (iv) follows from the fact that the sum is $(k-1)a>0$. For (v), $\Cor_R$ and
$B_R\le\hi$ give $(x_r-1)a\le B_R\le\hi$; the relation between the coefficients follows from (ii).
\end{proof}

\begin{proof}[Proof of Theorem~\ref{t:ind1}]
(i) In V, $h_\rV(t)=-j^2Q_{t_j}$ depends on one factor only, so its weighted average is $-j^2\sum_vw_vQ_v=-j^2a$
(because $\sum w=1$). Averaging over $j$ gives $-k\,a$. (ii) Since $P_\rV\subseteq P_\rL\subseteq P_\rLz$, the maximum
on each component increases; since the weights are nonnegative, so does the average. $Q\ge0$ gives $h\le0$.
(iii) Multiplication by a nonnegative number commutes with the maximum. (iv) On a component of positive weight all
factors have the same order, so the maximum equals $h_\rV$; components of weight $0$ do not enter the sum. (v) Take
places $u,v$ with $Q_u<Q_v$. The component $(u,v)$ for $j=1$ has weight $w_uw_v>0$, and there
$h_\rLz\ge-Q_u>-Q_v=h_\rV$; on the other components and labels one has $\ge$. For L use the component $(u,u,v)$ for
$j=2$.
\end{proof}

\begin{proof}[Proof of Proposition~\ref{p:rest}]
Each formula consists of a term of the form $\min_{i\in P}(s_i-c_i)$, with $0\le c_i\le c_{\max}$ and $c_{\max}$ a
local constant, and terms that do not involve the $q$-orders. Since
\[\min_is_i-c_{\max}\le\min_i(s_i-c_i)\le\min_is_i,\]
the remainder lies in an interval that does not depend on the $q$-orders. The readings without the hull and with the
large ball add to or subtract from the value with the hull quantities that do not involve the $q$-orders.
\end{proof}

\begin{proof}[Proof of Theorem~\ref{t:cheap}]
(i) If $N>C/((k-1)a)$, then $\Ht(D^N)$ fails, so by Theorem~\ref{t:cons} the two statements do not hold together.
(ii) If $\Theta_r>-k\,a$, then $g=\Theta_r+k\,a>0$, and $\IVb_R(D^N)$ gives $Ng\le C-B_R(D^N)\le C-\lo_R$, which fails
for large $N$. Hence $\Theta_r=-k\,a$, and substituting gives the equivalence. (iii) If $\Theta_r<-a$, then
$g=-a-\Theta_r>0$, and $\Cor_R(D^N)$ gives $Ng\le B_R(D^N)\le\hi_R$, which fails for large $N$. For the second half
insert $\Theta_r\ge-a$ into $\IVb_R(D^N)$. (iv) $M_R(D^N)/N=\Theta_r+B_R(D^N)/N$, and both ends of
$\lo_R/N\le B_R(D^N)/N\le\hi_R/N$ tend to $0$.
\end{proof}

\begin{proof}[Proof of Example~\ref{e:two}]
In L0 the maximum on a component is $-j^2Q_{\mathrm b}$ if all factors are at the bad place and $0$ otherwise. That
event has weight $w^{j+1}$, so $\Theta_\rLz=-wQ_{\mathrm b}A(w)$ and $a=wQ_{\mathrm b}$. $A$ is increasing in $w$, and
\[\sum_{j=1}^nj^2/2^j=6-(n^2+4n+6)/2^n\le n\]
gives $A(\frac12)\le1$. For V, Theorem~\ref{t:ind1}\,(i) and the remainder $0$ give $\IVb\iff0\le C$ and
$\Cor\iff(k-1)Na\le0$, which is false. IV in L0 follows from Theorem~\ref{t:cheap}\,(i) and $\Cor$.
\end{proof}

\section{Computations on explicit data}\label{s:data}

\subsection{Procedure}\label{ss:proc}
\tagP\ We use the seven initial $\Theta$-data of the catalogue~\cite{Cat}. The computation uses the program of the
walkthrough~\cite{WT}, version~1.1.0, extended by the hulls of twelve readings (\texttt{iut\_patterns.py}), their
assembly into the 36 readings (\texttt{iut\_readings36.py}) and the upper bound of IV (\texttt{iut\_height.py}). As
in~\cite{Gap} and~\cite{WT}, the set $\Vbad$ is that of the catalogue minus the places above $3$ and $5$, where the
ramification of $K$ depends on the construction of $F$ (Definition~3.1\,(b) of IUT~I allows such a subset; see~\cite[\S5]{Gap}). The places above
$2$, $3$, $5$ and $\ell$ therefore have theta value $1$, and so have those above the unfactored cofactor of $D_{2357}$.
In the readings with the hull these places only add nonnegative amounts to $M_R$; in the readings without the hull they
may also lower $M_R$, by at most their shell defect, which is not computed here. In the upper bound of IV they are
bounded by Step~(ii) and (R4), as in the computation of $U(D)$. In every reading the contribution of a prime at which the theta value is $1$ is at most the bound of IV there (with theta
value $1$ the lattice group gives the same as the units, and the iteration and the readings without the hull are at
most the hull of the container).

\subsection{The split point}\label{ss:x}
\begin{table}[ht]\centering\small
\caption{The ends and the split points ($\ell=7$ and $k=14/3$ except for $D_{2357}$). The last column is the number of
primes above which the $q$-orders differ at places of positive weight.}\label{tab:x}
\begin{tabular}{@{}lrrrrrrrr@{}}
\toprule
Datum & $\ell$ & $a$ & wall $(k-1)a$ & $C$ & $x_\rV$ & $x_\rL$ & $x_\rLz$ & primes\\
\midrule
$D_{211}$ & 7 & 21.801 & 79.937 & 1247.14 & 4.6667 & 1.7500 & 0.8750 & 5\\
$D_{109}$ & 7 & 19.114 & 70.085 & 1187.52 & 4.6667 & 1.7500 & 0.8750 & 6\\
$D_{2}$ & 7 & 17.149 & 62.881 & 1235.19 & 4.6667 & 2.2273 & 1.4955 & 2\\
$D_{3}$ & 7 & 80.589 & 295.49 & 2247.61 & 4.6667 & 2.3149 & 1.6094 & 15\\
$D_{18}$ & 7 & 43.925 & 161.06 & 1683.07 & 4.6667 & 1.7779 & 0.9113 & 15\\
$D_{15}$ & 7 & 36.862 & 135.16 & 1556.68 & 4.6667 & 1.7500 & 0.8750 & 15\\
$D_{2357}$ & 157 & 11.840 & 24463.9 & 16229.6 & 2067.17 & 0.1538 & 0.0769 & 9\\
\bottomrule
\end{tabular}
\end{table}

\tagP\ For all seven data, $x_\rV=k$ holds exactly as rational numbers ($\Theta_\rV+k\,a=0$ was checked with rational
coefficients of the $\log p$). For $D_{211}$, $D_{109}$ and $D_{15}$, above every bad prime there are exactly one bad and
one good place, each of weight $\frac12$. Hence $x_\rLz=A(\frac12)=\frac78$ and $x_\rL=2A(\frac12)=\frac74$: this is
exactly the situation of Example~\ref{e:two}.

\subsection{The verdicts in the 36 readings}\label{ss:verdicts}
\tagP\ For the six data that are completely factored (all except $D_{2357}$):
\begin{itemize}
\item The upper bound of IV holds in all 36 readings.
\item Corollary~3.12 holds for all six data in the container readings with the hull and in the readings of iteration
with the lattice group (18 readings).
\item In the readings of iteration with units (6 readings), Corollary~3.12 is short for all six data in V and L: the
counted places fall short, and it fails unless the places that are not counted (those above $2$, $3$, $5$ and $\ell$)
add at least the shortfall. In L0 only $D_3$ is short.
\item In the container readings without the hull (12 readings) the verdict is undecided, because the shell defect is
only bracketed. Only in L0 for $D_{109}$ does it hold, for the counted places.
\end{itemize}

\tagH\ That Corollary~3.12 and IV hold together on the actual data is because these data satisfy the height inequality
with a large margin (for the six data the wall $(k-1)a$ is 5--13\% of $C$). These verdicts test nothing. The effect of
the reading becomes visible in the family $q^N$.

\subsection{Replacing \texorpdfstring{$q$}{q} by \texorpdfstring{$q^N$}{q\^N}}\label{ss:qN}
\tagP\ The program multiplies only the $q$-orders by $N$ (the ramification indices and $\Vbad$ are kept, and
$\ell\nmid N$). For the six data and $N=2,3,25$ we checked that $a$ and the theta parts are multiplied by $N$ exactly,
and that $x$ and $C$ do not change. For $D_2$, $D_{211}$ and $D_{109}$ we checked $N=1,\dots,40$ (35 values, leaving out
the multiples of $\ell$). In every reading the range of the remainder $B_R$ is at most $19.9$, $25.9$ and $22.7$
respectively, about $3\%$ of the growth of $Na$ ($669$ to $850$). The largest range is already reached at $N=3$ and
does not grow beyond it. For $D_{109}$ the range $22.707$ equals $\Delta(D_{109})$, the value of~\cite{Gap}: at $N=3$
the numerators $n_u\in\{29,1\}$ of the $q$-orders, quadratic residues modulo $7$, become non-residues, so by
\cite[Proposition~3.6]{Gap} the label average of the rounding term of the lattice group goes from $\frac13$ to
$\frac23$ and the term doubles, from $\Delta(D_{109})$ to $2\Delta(D_{109})$ (in reading V with the container and the
hull the remainder goes from $228.556$ to $251.262$).

\begin{table}[ht]\centering\small
\caption{For which $N$ each statement fails in the family $q^N$. The remainder $B_R$ is that of the places counted in
Section~\ref{ss:proc}, replaced by its interval at $N=1$.}\label{tab:qN}
\begin{tabular}{@{}lll>{\raggedright\arraybackslash}p{4.3cm}l@{}}
\toprule
Datum & $\Ht$ fails & V: Corollary~3.12 & L0: Corollary~3.12 & L0: IV fails\\
\midrule
$D_{211}$ & $N>15.60$ & fails for $N>3.31$ & holds for every $N$ with the hull, and for $N\ge2$ in all twelve & $N>12.12$ to $15.15$\\
$D_{109}$ & $N>16.94$ & fails for $N>3.27$ & holds for every $N$ & $N>13.46$ to $16.34$\\
$D_{2}$ & $N>19.64$ & fails for $N>4.29$ & fails for $N>29.83$ & $N>18.05$ to $22.95$\\
$D_{3}$ & $N>7.61$ & fails for $N>3.14$ & fails for $N>17.67$ & $N>5.60$ to $9.30$\\
$D_{18}$ & $N>10.45$ & fails for $N>3.41$ & holds for every $N$ with the hull, and for $N\ge6$ in all twelve & $N>7.11$ to $10.33$\\
$D_{15}$ & $N>11.52$ & fails for $N>3.44$ & holds for every $N$ with the hull, and for $N\ge3$ in all twelve & $N>8.05$ to $11.22$\\
\bottomrule
\end{tabular}
\end{table}

A single number is the latest threshold over the twelve readings with the same (Ind1)-axis; a range runs from the least
to the largest threshold over these readings (\texttt{thresholds.py}). In the readings of L0 without the hull the
remainder has a negative lower end, so for $D_{211}$, $D_{18}$ and $D_{15}$ Corollary~3.12 is guaranteed only from the
stated $N$ on. In V, IV carries no load and holds for every $N$. For each reading the smaller of its two thresholds is at
most that of $\Ht$ (Theorem~\ref{t:cons}). Corollary~3.12 fails in V much earlier than $\Ht$ because its room $B_R$ is
much smaller than the whole room $C$ of IV: 38--71\% of $C$ is the bound at the primes $2$, $3$, $5$ and $\ell$
(Step~(ii) for the different and (R4) for the terms of $I^*$), which is not in $B_R$. The places above these primes change
$M_R$ by an amount that does not depend on $N$: in the readings with the hull it is nonnegative and at most that bound,
and without the hull it may also be negative, by at most their shell defect. Such an amount moves the thresholds of
Corollary~3.12 and of IV in opposite directions and leaves that of $\Ht$ unchanged.

\subsection{\texorpdfstring{$D_{2357}$}{D2357}}\label{ss:d2357}
\tagP\ For $D_{2357}$ ($\ell=157$) the counted places alone leave the height inequality short by $8234$
($(k-1)a=24464>C=16230$). In reading V this deficit appears on the side of Corollary~3.12 (IV holds); in L and L0 it
appears on the side of IV (Corollary~3.12 holds or is undecided, and the deficit of IV is $8362$ to $10308$ in the
readings with the hull and the small archimedean container, $8376$ to $10321$ with the large ball, and at least $7775$
without the hull, where the shell defect is only bracketed). In every reading one of the two statements has a deficit.

This is conditional. The places above the unfactored cofactor of 1564 digits are not counted. In the readings with the
hull they add nonnegative amounts to $M_R$, and they add to $U(D)$, where the estimate of Step~(ii) at $2$, $3$, $5$ and
$\ell$ also grows with the conductor; so either deficit may be filled. Moreover, $U(D)$ is the upper bound of Steps (v)--(vii); compared with the final uniform
form of IV (right-hand side $3.47\times10^9$), IV holds in every reading.

\section{Discussion}\label{s:disc}

\subsection{Cannot be avoided, cannot be made cheaper}
\tagH\ The question had two parts: whether choosing the readings differently avoids the height inequality
(avoidance), and whether it makes the proof cheaper (cost).
\begin{itemize}
\item \emph{It cannot be avoided.} In every reading, the two statements together are the same height inequality
(Theorem~\ref{t:cons}). This is because the ends take no reading, which rests on the wording of Corollary~3.12 and on
Step~(v) of IV (Section~\ref{s:setting}).
\item \emph{It cannot be made cheaper.} A reading only moves the content of the wall between the two statements; the
total $(k-1)a$ does not change (Theorem~\ref{t:split}\,(ii)). No reading makes both loads nonpositive
(Theorem~\ref{t:split}\,(iv)). Making one statement hold for every $N$ makes the other heavier by the same amount, so
that it carries at least the whole wall (Theorem~\ref{t:cheap}\,(ii), (iii)). Rebuilding IV for each reading into a form
that holds automatically weakens the resulting inequality by exactly that amount (Theorem~\ref{t:split}\,(v)).
\end{itemize}

\tagH\ If IUT makes a contribution of its own to the height inequality, it should lie in computing the two sides by a
method different from the elementary one. But in each of the 36 readings the two sides are computed by the same single
procedure (Section~\ref{s:setting}), and the numbers at the ends agree with the elementary computation. What depends on
the reading is only the middle term $-|\log\TT|$, and it cancels when the two statements are combined.

\subsection{Two typical readings}\label{ss:types}
\tagH\ In reading V the load of IV is $0$, and Corollary~3.12 carries the whole wall: Corollary~3.12 is itself the
height inequality $(k-1)a\le B_R$. This corresponds to the side of the dilemma of Scholze and Stix~\cite[\S2.2]{SS} on
which the factor $j^2$ is not absorbed by any blurring: Corollary~3.12 is then a height inequality that has to be
proved as it stands, and in the family $q^N$ it fails for large $N$. The inequality $(k-1)Q\le E$ of Part~I of~\cite{II} has this form.

\tagH\ In reading L0, for data whose good and bad places lie above the same prime (Belyi points), one has $x\le1$ and
Corollary~3.12 carries no positive load: it holds whatever $q$ is, and IV carries at least the whole wall. This is the
situation of the option of Scholze and Stix labelled (c) in~\cite[Proposition~5.3 and \S15.1]{II}, in which
\q{the blurring must be by a factor of at least $\mathrm{O}(\ell^2)$ rendering the inequality thus obtained useless}
\cite[\S2.2]{SS}; this sentence has been formalized in Lean~4, with the constant $\ell(\ell+1)/12$, our $k$, in place
of $\mathrm{O}(\ell^2)$~\cite{sfingali}. The statement of lana that the literal reading does not suffice for ABC
(Section~\ref{ss:prior}) can be read as a description of this side from the standpoint of IV.

\tagH\ The other readings put the split point $x$ between these two. The point $x$ is decided by (Ind1) alone, and for
data whose $q$-orders agree, every reading has $x=k$, as V does.

\subsection{What could change the conclusion}\label{ss:change}
\tagH\ The conclusion of this note changes only if one of the following happens.
\begin{itemize}
\item The ends depend on the reading. For the left end this would contradict the wording of Corollary~3.12 (the
log-volume $-|\log\qq|$ is not subject to the indeterminacies). \tagX\ It is the position of Joshi, for whom
\q{the Tate parameter changes in non-trivial ways with changes of Arithmetic Holomorphic Structures and in particular
with $\mathrm{Ind}_1$, $\mathrm{Ind}_2$ and $\mathrm{Ind}_3$ (to be clear Mochizuki does not prove this--that is why the
theory is incomplete)}~\cite[(8)\,(iv)]{Joshi}; on that view the conservation of this note does not apply.
\item The theta part is smaller than $-k\,a$. The possible images include the $\Theta$-pilot object itself (the identity
permutation and the identity automorphism), so the value of the hull is at least that of the $\Theta$-pilot. A Kummer
distortion proportional to $a$ that acts in the negative direction, however, could give $x>k$. Even then the sum of the
loads does not change (the proof of Theorem~\ref{t:split}\,(ii) does not use $\Theta\ge-k\,a$), and the load moves to
the side of Corollary~3.12.
\item A reading outside the 36 has a theta part that is not proportional to the $q$-orders. Then there is no slope
(Theorem~\ref{t:cheap}\,(iv)), and the split of this note cannot be used.
\end{itemize}

\section{What is not shown}\label{s:notshown}
\begin{itemize}
\item That the 36 readings are all the readings of the original papers.
\item A Lean proof that each axis changes only the remainder (Proposition~\ref{p:rest}). Lean uses as a hypothesis that
the remainders lie in intervals that do not depend on $N$, and defines the theta part as the maximum over the set of
positions. That the axes other than (Ind1) do not change the theta part was checked by a paper argument from the
formula for each axis and by the programs.
\item That the family $q^N$ is a family of actual curves.
\item Which reading the original papers intend, and whether Corollary~3.12 holds.
\item The verdict for $D_{2357}$, which depends on the unfactored cofactor.
\end{itemize}

\section{Checks and reproducibility}\label{s:checks}
The folder \texttt{reading-family/} of the repository
\url{https://github.com/turahigashi/iut-explicit-initial-theta} (tag \texttt{v0.24.0}), also deposited with this note,
contains the Lean files, the programs and their outputs.
\begin{itemize}
\item \emph{Lean.} \texttt{lean/ReadingFamily.lean} (1024 lines) was checked with Lean~4.30.0 and Mathlib at tag
\texttt{v4.30.0}. The file \texttt{lean/ReadingFamilyAxiomCheck.lean} is the same file followed by \lean{\#print axioms}
for all 96 named declarations. Every declaration depends at most on \lean{propext}, \lean{Classical.choice} and
\lean{Quot.sound}; the definition of the readings, \lean{Reading.Admissible}, depends on none. There is no
\lean{sorry}, no \lean{axiom}, no \lean{opaque} and no \lean{native\_decide}; there are no errors and no warnings. As a
witness that the hypotheses are not empty, the two-place family \lean{belyiFamily} is constructed.
\item \emph{Programs} (Python~3.10 or later with sympy; their comments and output are in Japanese, as for the program
of~\cite{WT}). \texttt{scripts/iut\_compute.py} is the program of~\cite{WT}, version~1.1.0, unchanged.
\texttt{iut\_patterns.py} computes the twelve readings with the hull; with \texttt{-{}-check} it confirms that its
container--units readings agree with the original program, that its difference of the two (Ind2) models in reading V
is $\Delta(D)$ of~\cite{Gap}, that its closed formulas agree with the enumeration of the tuples, and that the intervals
for $\kappa$ contain the exact values in model fields. Its output is identical, character for character, whether it
uses version 1.0.0 or version 1.1.0 of the program of~\cite{WT}. \texttt{iut\_readings36.py} assembles the 36 readings, checks
$\Theta_\rV=-k\,a$ exactly, and checks numerically the order $\rV\le\rL\le\rLz\le0$ of the theta parts and the sums of
the loads and of the rooms, for all data; with \texttt{-{}-scale} it checks the family $q^N$. \texttt{scan\_scale.py}
scans $N=1,\dots,40$, \texttt{thresholds.py} computes Table~\ref{tab:qN}, \texttt{iut\_height.py} computes $U(D)$, and
\texttt{elementary\_height.py} and \texttt{compare\_elementary\_iut.py} give the comparison with the elementary
computation of Section~\ref{s:setting}.
\end{itemize}

\section*{Acknowledgements and use of AI}
The question studied in this note---whether choosing the readings differently avoids the height inequality or makes the
proof cheaper---arose in the author's discussions with generative AI tools; the author judged it important and decided
to make it the subject of this note. Its ultimate aim, to understand exactly the algorithm by which IUT computes the
height inequality, is the author's research direction. Generative AI tools (large language models) were used
extensively, under the author's direction, in finding the proofs, in the computations, in the verification in Lean and
in the writing. The author takes full responsibility for the contents of this note.

\begin{thebibliography}{sfingali}\raggedright
\bibitem[Cat]{Cat}
Toshihisa Urahigashi, \emph{Explicit initial $\Theta$-data and the two readings of Step (v) of IUT IV, Theorem 1.10:
a catalogue}, version 1.0.0, Zenodo, 2026, \href{https://doi.org/10.5281/zenodo.23870448}{doi:10.5281/zenodo.23870448}.

\bibitem[Fifth]{Fifth}
Toshihisa Urahigashi, \emph{Two constants in the numerical estimates of inter-universal {T}eichm\"uller theory {IV}:
$12/(\ell(\ell+1))$ and $12/\ell^2$}, fifth report, version 1.0.0, Zenodo, 2026,
\href{https://doi.org/10.5281/zenodo.22659796}{doi:10.5281/zenodo.22659796} (all versions).

\bibitem[Form]{Form}
Shinichi Mochizuki, \emph{On the formalization of IUT: a preliminary progress report} (joint work in progress with
Y.~Hoshi, G.~Yamashita, Y.~Yang), slides, April 2026, 14 pp.

\bibitem[Fourth]{Fourth}
Toshihisa Urahigashi, \emph{Split bad primes and the tuple structure of inter-universal {T}eichm\"uller theory}, fourth
report, version 1.0.4, Zenodo, 2026, \href{https://doi.org/10.5281/zenodo.22648944}{doi:10.5281/zenodo.22648944} (all
versions).

\bibitem[Gap]{Gap}
Toshihisa Urahigashi, \emph{The gap between the two walls is positive: a closed form for the (Ind2)-terms of
lana-agents/iut}, version 1.0.0, Zenodo, 2026, \href{https://doi.org/10.5281/zenodo.23954194}{doi:10.5281/zenodo.23954194}.

\bibitem[II]{II}
Toshihisa Urahigashi, \emph{Reading the IUT dispute through indexed identifications: an educational model in
undergraduate mathematics}, version 2.0.0, Zenodo, 2026,
\href{https://doi.org/10.5281/zenodo.23933108}{doi:10.5281/zenodo.23933108}.

\bibitem[IUT3]{IUT3}
Shinichi Mochizuki, \emph{Inter-universal {T}eichm\"uller theory {III}: canonical splittings of the log-theta-lattice},
Publ. Res. Inst. Math. Sci. \textbf{57} (2021), 403--626.

\bibitem[IUT4]{IUT4}
Shinichi Mochizuki, \emph{Inter-universal {T}eichm\"uller theory {IV}: log-volume computations and set-theoretic
foundations}, Publ. Res. Inst. Math. Sci. \textbf{57} (2021), 627--723.

\bibitem[Joshi]{Joshi}
Kirti Joshi, \emph{Comments on the LANA Project Report of Kato et al.}, 30 July 2026,
\url{https://sites.arizona.edu/kirti-joshi/files/2026/07/Comments-on-the-LANA-Project-Report-of-Kato-et-al.pdf}.

\bibitem[lana]{lana}
lana-agents, \emph{iut}: a Lean~4 formalization of inter-universal {T}eichm\"uller theory, GitHub repository,
\url{https://github.com/lana-agents/iut}, commit \texttt{553076a} (3 November 2026).

\bibitem[LANA]{LANA}
Project LANA, \emph{Project LANA interim report on IUT theory}, 20 July 2026,
\url{https://ncatlab.org/nlab/files/LANAProject-Report-July2026.pdf}.

\bibitem[MT]{MT}
A.~Minamide and S.~Tsujimura, \emph{Anabelian geometry for {H}enselian discrete valuation fields with quasi-finite
residues}, RIMS-1973, Research Institute for Mathematical Sciences, Kyoto University, June 2023.

\bibitem[Richards]{Richards}
David Richards, \emph{motizukipdf---{M}ochizuki's papers, plus a toolkit for the data types of {IUT}}, GitHub
repository \url{https://github.com/oneman/motizukipdf}, consulted at commit \texttt{2665db1} (30 October 2026); Lean port
\lean{lean/IUT/LeanForm.lean}.

\bibitem[sfingali]{sfingali}
sfingali, \emph{The {S}cholze--{S}tix objection, formally: the arithmetic is agreed, the $q$-side is absent}, Zenodo,
2026, \href{https://doi.org/10.5281/zenodo.21757640}{doi:10.5281/zenodo.21757640}.

\bibitem[SS]{SS}
Peter Scholze and Jakob Stix, \emph{Why abc is still a conjecture}, 2018.

\bibitem[TK]{TK}
Takkun-kohinata, \emph{IUT\_LEAN}, Lean repository, \url{https://github.com/Takkun-kohinata/IUT_LEAN}, commit
\texttt{9d56c46} (15 July 2026); file \lean{IutLean/Multiradial/ThetaHullDichotomy.lean}.

\bibitem[W311]{W311}
Toshihisa Urahigashi, \emph{Where Theorem 3.11 sits: the log-volume content of the results cited in the proof of IUT
III, Corollary 3.12, does not imply its conclusion}, Zenodo, 2026,
\href{https://doi.org/10.5281/zenodo.24028832}{doi:10.5281/zenodo.24028832} (all versions).

\bibitem[WT]{WT}
Toshihisa Urahigashi, \emph{A programmer's walkthrough of inter-universal {T}eichm\"uller theory: running the
multiradial algorithm on explicit initial $\Theta$-data}, version 1.1.0, Zenodo, 2026,
\href{https://doi.org/10.5281/zenodo.23955994}{doi:10.5281/zenodo.23955994}.

\bibitem[Zeng]{Zeng}
Chunping Zeng, \emph{An integral-structure inflation in the hull operation of {IUT} {C}orollary 3.12}, Zenodo, 2026,
\href{https://doi.org/10.5281/zenodo.23811410}{doi:10.5281/zenodo.23811410}.
\end{thebibliography}

\end{document}
